The analysis stage converts the three benchmark artifact layers into one canonical table, then derives metrics, aggregate tables and plots from that single source. The detailed implementation walk-through is preserved in Appendix~\ref{app:extended-pipeline}; this section keeps only the information needed to understand the evaluation flow.

\begin{center}\small
\begin{tabular}{@{}ccccc@{}}
\toprule
\textbf{3} & \textbf{1} & \textbf{8} & \textbf{21} & \textbf{7}\\
artifact layers & canonical DataFrame & aggregate tables & registered plots & capability profiles\\
\bottomrule
\end{tabular}
\end{center}

\subsubsection{The pipeline at a glance}

The pipeline has a fixed data path:

\begin{center}
\small
\texttt{raw/} $\rightarrow$ \texttt{parsed/} $\rightarrow$ \texttt{scored/} $\rightarrow$ \texttt{df\_metrics} $\rightarrow$ \texttt{aggregate tables} $\rightarrow$ \texttt{plots + report.json}
\end{center}
The analysis script joins these layers with the selected PDDL task files, so every row can be evaluated against the exact domain and problem that generated it.

\begin{table}[H]\small
\centering
\begin{tabularx}{\linewidth}{@{}l Y@{}}
\toprule
\textbf{Stage} & \textbf{Purpose}\\
\midrule
Load artifacts & Creates one row per tested instance, preserving model, domain, difficulty, protocol, validity and iteration information.\\
\addlinespace
Index PDDL context & Extracts legal actions, objects, initial state and numeric fluents from the task files.\\
\addlinespace
Compute row metrics & Adds hallucination, executability, precondition-awareness and CoT-alignment metrics to each row.\\
\addlinespace
Aggregate & Builds model-level, domain-level and ranking tables from row-level metrics.\\
\addlinespace
Visualize and export & Generates the registered plots and serializes the report payload.\\
\bottomrule
\end{tabularx}
\caption{Compact view of the analysis pipeline.}
\label{tab:pipeline-compact}
\end{table}

\subsubsection{The canonical table: \texorpdfstring{\texttt{df\_metrics}}{df\_metrics}}

\texttt{df\_metrics} is the central object of the analysis. It starts from the artifact rows and is horizontally extended with computed columns, its grain is one benchmark instance not one retry attempt.\\
This is important because success rates, hallucination rates and executability scores are averaged over comparable units.

The retry structure is still available: the loader preserves per-attempt lists in private columns and only the metrics that need them, such as FASR, IWSR, iteration profiles and CoT alignment, read those lists.

\begin{table}[H]\small
\centering
\begin{tabularx}{\linewidth}{@{}l Y@{}}
\toprule
\textbf{Input} & \textbf{Role}\\
\midrule
Artifact rows & Provide identity fields and direct outcomes.\\
\addlinespace
Domain context & Provides action schemas, legal action names, predicates and numeric fluents.\\
\addlinespace
Problem context & Provides objects, initial propositions and numeric initial values for the specific instance.\\
\addlinespace
Warnings & Records missing or malformed inputs without dropping rows.\\
\bottomrule
\end{tabularx}
\caption{Inputs used to build the canonical metric table.}
\label{tab:dfmetrics-inputs-compact}
\end{table}

\subsubsection{\texorpdfstring{\texttt{PDDLSimulator}}{PDDLSimulator} role}

Some metrics require more than the final VAL decision. A plan can be invalid because it fails immediately, because it follows legal actions in the wrong order or because it invents a state that never existed; to separate these cases, the analysis uses a lightweight forward simulator over the parsed PDDL fragments.

The simulator starts from the problem initial state, applies each known action step by step and stops at the first unsatisfied precondition. If a missing proposition was true before, the error is classified as sequencing; if it was never true, it is classified as state fabrication.\\
VAL remains the official correctness oracle: the simulator is a diagnostic layer for explaining invalid plans.